\section*{Erd\H{o}s Problem 784: the reciprocal budget in a small sieve}

\paragraph{Statement and claim.}
For \(A\subseteq[1,x]\cap\mathbb N\), assume
\(\sum_{a\in A}1/a\leq C\), and count integers up to \(x\)
divisible by no element of \(A\).
The answer depends on both \(C\) and whether \(1\) is allowed.
Literally, the proposed conclusion holds for \(0<C<1\) and
fails for \(C\geq1\).
With the intended additional condition \(1\notin A\), it holds
for \(0<C\leq1\) and fails for every \(C>1\).
We give the elementary part and derive the complete classification
from an explicitly stated established sieve theorem.
See \url{https://www.erdosproblems.com/784}.

\paragraph{The literal statement.}
Write \(X=\lfloor x\rfloor\). The union bound gives
\[
 \#\{m\leq X:\exists a\in A,\ a\mid m\}
 \leq\sum_{a\in A}\lfloor X/a\rfloor\leq CX.
\]
For \(0<C<1\), at least \((1-C)X\gg_C x\) integers survive,
which is stronger than any requested fixed logarithmic loss.
For \(C\geq1\), the choice \(A=\{1\}\) has admissible reciprocal
sum and leaves no survivors. Thus the literal statement
cannot include \(C=1\).

\paragraph{The intended convention \(1\notin A\).}
Define \(H_C(x)\) as the minimum survivor count over
\(A\subseteq\{2,\ldots,\lfloor x\rfloor\}\) with reciprocal
sum at most \(C\).
The external input is Weingartner's Theorem 5, incorporating
and sharpening results of Ruzsa and Saias:
for every fixed \(Z>1\), uniformly for \(1\leq C\leq Z\),
\[
 H_C(x)\asymp_Z\frac{x^{\exp(1-C)}}{\log x}.                    \tag{1}
\]
Source: A. Weingartner, \emph{The Schinzel--Szekeres function},
Research in Number Theory 11 (2025), Paper 63,
\url{https://arxiv.org/abs/2310.13038}.
The theorem uses exactly the reciprocal constraint and
excludes \(1\). Its minimum may equivalently range over
larger sets, because elements exceeding \(x\) affect neither
the survivor count nor the usefulness of the budget.
The substantial sieve estimates behind (1) are external,
not reproved here.

At \(C=1\), (1) gives \(H_1(x)\asymp x/\log x\), so \(c=1\)
works. For \(C<1\), the earlier union bound still applies.
For fixed \(C>1\), set \(\theta=\exp(1-C)<1\).
The upper bound in (1) supplies admissible sets with
at most \(K_Cx^\theta/\log x\) survivors.
For every fixed real \(c>0\),
\[
 \frac{K_Cx^\theta/\log x}{x/(\log x)^c}
 =K_Cx^{-(1-\theta)}(\log x)^{c-1}\longrightarrow0.
\]
Consequently no exponent depending only on this \(C\)
can give the proposed universal lower bound.
The minimizing set may depend on \(x\); this is allowed
by the quantifiers in the question.

\paragraph{Why a prime-only intuition is insufficient.}
For a finite set of primes \(P\), the density of integers
not divisible by any of them is
\(\prod_{p\in P}(1-1/p)\).
If \(\sum_{p\in P}1/p\leq C\), then
\[
 \prod_{p\in P}(1-1/p)
 \geq\exp\left(-2\sum_{p\in P}1/p\right)\geq e^{-2C},
\]
using \(-\log(1-u)\leq2u\) for \(0\leq u\leq1/2\).
This calculation concerns a fixed periodic sieve. It is
not by itself a uniform finite-\(x\) bound when the prime
set varies with \(x\). It nevertheless identifies the
independence property that one must not impose on the
composite-divisor sets in (1).

\paragraph{Scope.}
Both conventions are completely classified relative to the
stated Weingartner theorem. The elementary union bound,
the obstruction from \(1\), and the comparison of powers
with logarithms are proved directly.

\paragraph{Completion estimate.}
\textbf{COMPLETION: 100\%}
